\chapter{Auctions and Algorithms}
\label{ch:auctions_and_algorithms}

This chapter sets up the auction environment shared by all experiments and describes the three families of bidding algorithms that I study. The environment is a standard repeated single-item auction with a discrete bid grid, and the algorithms are Q-learning, contextual bandits, and budget-constrained pacing. Each algorithm is introduced first in plain-language economic terms, before the formal update rule is written down. Equilibrium benchmarks are stated briefly, and their derivations appear in Appendix~\ref{ch:equilibria}. A reader who wants the intuition without the mathematics can skim the opening paragraph of each algorithm section and skip the equations.

\section{Auction Environment}
\label{sec:env}

The environment is a sequence of single-item auctions indexed by $t = 1, 2, \ldots, T$, with $n$ bidders competing in each round. In round~$t$, bidder~$i$ submits a bid $b_{it}$ from a finite grid $\mathcal{B} = \{0, 0.1, 0.2, \ldots, 1.0\}$ and pays a price determined by the auction rule. The grid has eleven equally spaced points, which is fine enough to let distinct bidding behaviours emerge yet coarse enough that Q-learning agents visit each action many times over the course of a run. Robustness checks in Appendix~\ref{ch:robustness} show that factor rankings are stable when the grid is widened to six levels or narrowed to twenty-one levels, so the central findings are not artefacts of this particular choice.

A reserve price $r \geq 0$ filters out bids below the threshold. If every submitted bid is below $r$, no sale occurs in that round and payoffs are zero for all bidders. Otherwise, any bid strictly below $r$ is excluded from contention, and ties for the top bid are broken uniformly at random among tying bidders.

Two payment rules are studied. Under the first-price rule, the winner pays her own bid; under the second-price rule, the winner pays the highest bid among the other valid bids (or the reserve price if only one valid bid remains). Writing $p_{it}$ for the payment bidder~$i$ makes in round~$t$, her per-round payoff is
\[
u_{it} \;=\; (v_{it} - p_{it}) \cdot \mathbb{1}[\text{bidder } i \text{ wins round } t],
\]
where $v_{it}$ is the bidder's valuation, drawn as described below. Losers pay nothing and receive zero payoff.

\subsection{Valuations}

Experiment~1a uses constant valuations $v_{it} = 1$ for all bidders and all rounds. This is the simplest possible value model and serves as a baseline free from signal heterogeneity. Experiments~1b, 2a, and 2b introduce affiliated signals through a linear-affiliation function
\[
v_{it}(s_{it}, s_{-i,t}) \;=\; (1 - 0.5\,\eta)\,s_{it} \;+\; \frac{0.5\,\eta}{n-1}\sum_{j \neq i} s_{jt},
\]
where $s_{it}$ is bidder~$i$'s private signal drawn independently from a uniform distribution on $\{0, 0.1, \ldots, 1.0\}$, and $\eta \in \{0, 0.5, 1\}$ governs the strength of affiliation. When $\eta = 0$, valuations collapse to independent private values; when $\eta = 1$, they lie closer to a common-value setting. The coefficients are chosen so that the expected valuation $\mathbb{E}[v_{it}]$ is the same at all affiliation levels, which isolates the effect of affiliation from any mean shift.

Experiments~3a and~3b replace the discrete-signal model with log-normal valuations to match the stylised setting of the autobidding literature. Each bidder draws an asymmetry mean $\mu_i$ uniformly from $[0.5, 1.5]$ once per seed, and in each round the valuation is $v_{it} \sim \mathrm{LogNormal}(\mu_i, \sigma)$. The dispersion parameter $\sigma \in \{0.3, 0.7\}$ controls how heterogeneous valuations are across rounds, which in turn affects the rate at which budget constraints bind.

\subsection{Repeated Structure and State}

A defining feature of the environment is that bidders may condition future bids on information from prior rounds. Depending on the experiment, a bidder's state may include the previous round's winning bid, her own signal, or a context vector of recent outcomes. The discount factor $\gamma \in [0, 1]$ governs how strongly a bidder weighs future rewards relative to present rewards. Agents with $\gamma$ close to 1 behave patiently and optimise long-run payoffs; agents with $\gamma$ closer to 0 behave myopically and maximise immediate payoff.

Every experiment repeats for a large number of rounds (typically $100{,}000$ under Q-learning and bandits; $100 \times 1{,}000$ episode-round cycles under pacing) so that any long-run behaviour has time to emerge.

\section{Q-Learning}
\label{sec:qlearning}

Q-learning is one of the oldest and simplest reinforcement-learning rules. The agent holds a lookup table whose rows are the states it can observe and whose columns are the actions it can take. Each entry $Q(s, a)$ is a running estimate of the long-run reward the agent expects if it is in state~$s$, takes action~$a$, and behaves optimally thereafter. After each round, the agent observes the actual reward and the next state, and shifts the corresponding entry towards a target that combines the reward just earned with its own current estimate of the best continuation value. In a stationary environment (one whose reward distribution does not drift over time) with a single agent and enough random exploration, the table converges to the optimal action-value function from classical dynamic programming. I note this as standard background only. The classical convergence result is proved for a single agent learning against a fixed environment; in my experiments, every bidder is a Q-learner, so the environment each agent faces is non-stationary. I therefore do not claim that my Q-learning agents reach any particular fixed point. The empirical question I ask is whether their long-run revenue outcomes respond systematically to changes in market structure, and the factorial analysis in Chapter~\ref{ch:qlearning} answers that. In the thesis experiments, the states are discretised records of recent auction outcomes (such as the previous winning bid), the actions are bid levels on the eleven-point grid, and the rewards are auction payoffs. The agent therefore learns to bid by trial and error, without ever computing an equilibrium strategy in closed form.

Formally, the standard asynchronous update is
\begin{equation}
Q(s_t, a_t) \;\leftarrow\; Q(s_t, a_t) \;+\; \alpha \left[ r_t + \gamma \max_{a'} Q(s_{t+1}, a') - Q(s_t, a_t) \right],
\label{eq:qlearning_async}
\end{equation}
where $\alpha \in (0, 1]$ is the learning rate and $r_t$ is the immediate reward earned in round~$t$.\footnote{The symbol $r$ is used with two different meanings in this chapter. In \Cref{sec:env} the lowercase $r$ denotes the reserve price, a non-negative threshold below which bids are excluded. In this section the subscripted $r_t$ denotes the reward a Q-learning agent receives in round~$t$. The two meanings are standard in their respective literatures, so I keep them, but the reader should disambiguate by looking for the subscript.} Only $Q(s_t, a_t)$ changes in any given update, and other entries are untouched. A synchronous variant updates $Q(s_t, a)$ for every possible action~$a$ simultaneously using counterfactual rewards, meaning the rewards the agent would have received had it taken each alternative action instead. Computing counterfactual rewards is possible here because the simulator knows the true valuations and can compute the payoff the agent would have earned under any bid. This variant accelerates convergence at the cost of extra computation. Both synchronous and asynchronous updates are studied in Experiment~1a.

Agents must balance exploration (trying actions whose reward is still uncertain) and exploitation (playing the action with the highest current estimate). Two standard rules govern action selection. Under $\varepsilon$-greedy exploration, the agent picks $\arg\max_a Q(s, a)$ with probability $1 - \varepsilon$ and picks uniformly at random with probability $\varepsilon$. Under Boltzmann exploration (also called softmax because of its functional form), the agent samples action~$a$ with probability proportional to $\exp(\beta \, Q(s, a))$, which smoothly increases the chance of high-value actions without ever eliminating the chance of low-value ones. Higher $\beta$ concentrates the distribution on the high-value actions and makes the rule closer to greedy play; lower $\beta$ makes the distribution more uniform and increases exploration. In both cases the exploration parameter decays over the horizon, shifting the agent from random probing early in training to near-greedy play late in training.\footnote{The $\varepsilon$ (or inverse Boltzmann temperature) starts at $1.0$, decays over the first 90\% of episodes to a floor of $0.01$, and then switches to pure exploitation for the final 10\%. Linear and exponential decay schedules are both studied as factors in Experiment~1a.}

\section{Contextual Bandits}
\label{sec:bandits}

A contextual bandit is a simpler cousin of Q-learning that omits state transitions. Before each decision, the agent receives a context vector $\mathbf{x}_t$ (a list of real-valued features that summarise whatever is observable at round~$t$, such as the bidder's own signal and the recent winning bid). Given that context, the agent picks the action it believes will yield the highest immediate reward and then updates its beliefs about how rewards depend on context. Because there are no forward-looking state dynamics, the discount factor plays no role. Contextual bandits are the algorithmic workhorse of online recommendation systems and display advertising, where the platform observes user features and must select a treatment quickly without tracking long-run dynamics. I study two variants that differ in how they balance exploration (trying actions whose reward is uncertain) against exploitation (using the best action known so far).

The first variant is LinUCB, which stands for linear upper confidence bound. The agent assumes that the expected reward for action~$a$ is a linear function of the context, $\mathbb{E}[r \mid \mathbf{x}, a] \approx \boldsymbol{\theta}_a^{\top} \mathbf{x}$, estimates the coefficient vector $\boldsymbol{\theta}_a$ by ordinary least squares on the observed context-reward pairs, and picks the action with the highest upper confidence bound on its estimated reward. Concretely, the agent maintains a matrix $\mathbf{A}_a = \lambda \mathbf{I} + \sum_{\tau} \mathbf{x}_{\tau} \mathbf{x}_{\tau}^{\top}$ and a vector $\mathbf{b}_a = \sum_{\tau} r_{\tau} \mathbf{x}_{\tau}$, where $\lambda > 0$ is a regularisation constant that keeps the matrix invertible when data are scarce. It then computes $\hat{\boldsymbol{\theta}}_a = \mathbf{A}_a^{-1} \mathbf{b}_a$ and picks
\begin{equation}
a_t \;=\; \arg\max_{a} \left[ \hat{\boldsymbol{\theta}}_a^{\top} \mathbf{x}_t \;+\; c \, \sqrt{\mathbf{x}_t^{\top} \mathbf{A}_a^{-1} \mathbf{x}_t} \right],
\label{eq:linucb}
\end{equation}
where $c > 0$ is an exploration parameter that scales the confidence bonus. The first term is the expected reward under current beliefs; the second term is an optimism bonus that encourages probing actions whose reward is poorly estimated.

The second variant is Thompson Sampling, named after the statistician W.\ R.\ Thompson, which takes a Bayesian view. The agent maintains a posterior distribution over $\boldsymbol{\theta}_a$ (a probability distribution that represents its beliefs about the unknown coefficient vector after seeing the data so far). At each round it draws a random sample $\tilde{\boldsymbol{\theta}}_a$ from the posterior and picks the action that would be best if the sample were the truth,
\begin{equation}
a_t \;=\; \arg\max_{a} \; \tilde{\boldsymbol{\theta}}_a^{\top} \mathbf{x}_t, \quad \tilde{\boldsymbol{\theta}}_a \sim \mathcal{N}\!\left(\hat{\boldsymbol{\theta}}_a,\, \sigma^2 \mathbf{A}_a^{-1}\right).
\label{eq:thompson}
\end{equation}
The variance parameter $\sigma^2$ plays a role analogous to $c$ in LinUCB. Larger $\sigma^2$ produces wider posterior draws and more exploration.\footnote{The same letter $\sigma$ appears elsewhere in the thesis for the log-normal dispersion parameter in the pacing experiments (\Cref{sec:pacing}). The two roles are distinct, and I retain the notation from the respective literatures to avoid renaming symbols that the reader may recognise from the source papers.}

Experiment~2a also introduces an optional memory-decay parameter $\delta \in (0, 1]$ that exponentially discounts historical observations. With $\delta < 1$, the effective observation window shrinks to roughly $1 / (1 - \delta)$ rounds, preventing the agent from locking into the dominant action based on early experience. This mechanism tests whether the rigidity of full-memory bandit learners contributes to bid suppression, following the convergence argument of \citet{Douglas2024}.

\section{Budget-Constrained Pacing}
\label{sec:pacing}

Pacing algorithms solve a different problem from Q-learning and bandits. Instead of learning which action is best in the abstract, a pacing agent is given a budget and must spread its spending across a horizon of auctions so that the budget lasts until the end. The agent has a budget $B^i$ and $T$ rounds. Each round delivers an exogenous valuation $v_t^i$, and the agent must decide how aggressively to bid so that the total spend approximately equals the budget by the final round. The canonical approach, which I adopt, shades each bid by a multiplier that stands in for the shadow price of the budget constraint from a Lagrangian relaxation of the agent's optimisation problem, familiar to economists as the dual variable of a constrained optimisation.

Experiment~3a uses multiplicative dual pacing. The agent maintains a dual variable $\mu_t^i$ and bids
\begin{align}
b_t^i \;&=\; \min\!\Big(\frac{v_t^i}{\mu_t^i},\; B^i - S_t^i\Big) & &\text{(value-maximiser),} \label{eq:bid_vmax}\\
b_t^i \;&=\; \min\!\Big(\frac{v_t^i}{1 + \mu_t^i},\; B^i - S_t^i\Big) & &\text{(utility-maximiser),} \label{eq:bid_umax}
\end{align}
where $S_t^i$ is the cumulative spend so far. After each round, the dual variable updates via an exponential rule
\begin{equation}
\mu_{t+1}^i \;=\; \operatorname{clip}\!\Big(\mu_t^i \cdot \exp\!\big(\alpha_p \, (c_t^i - B^i / T)\big),\; \mu_{\min},\; \mu_{\max}\Big),
\label{eq:ac_dual_update}
\end{equation}
where $\alpha_p = 1/\sqrt{T}$ is the dual step size and $c_t^i$ is the payment actually made in round~$t$. The operator $\operatorname{clip}(x, a, b)$ simply clamps $x$ to the interval $[a, b]$, setting it to $a$ if $x < a$ and to $b$ if $x > b$. Here it prevents the dual variable from drifting to extreme values that would cause numerical problems. The bounds $[10^{-4}, 100]$ are chosen loose enough that they rarely bind in the simulation output, and post-hoc checks confirm that widening them does not change the factor rankings reported in Chapter~\ref{ch:pacing}.\footnote{The dual-update rule follows \citet{Balseiro2019Learning}, who proves individual-agent regret guarantees of order $\sqrt{T}$ for this scheme.} When per-round spend exceeds the target spend rate $B^i / T$, the dual variable rises and shades the bid down, and underspending reverses the adjustment.

Experiment~3b replaces the multiplicative dual update with a proportional-integral (PI) feedback controller. The controller tracks the cumulative spending error
\begin{equation}
e_t^i \;=\; \frac{t}{T}\,B^i \;-\; \sum_{\tau \leq t} c_{\tau}^i,
\label{eq:pi_error}
\end{equation}
and updates its shading multiplier by
\begin{equation}
\lambda_{t+1}^i \;=\; \operatorname{clip}\!\Big(\lambda_t^i + K_P \, e_t^i + K_I \sum_{\tau} e_{\tau}^i,\; 0.01,\; 1.5\Big),
\label{eq:pi_update}
\end{equation}
with proportional and integral gains $K_P = 0.30 \times \text{aggressiveness}$ and $K_I = 0.05 \times \text{aggressiveness}$. A positive error means the agent is underspending relative to the linear budget trajectory, which raises $\lambda$ and encourages higher bids. The derivative term is omitted, following common practice in stochastic pacing implementations where differentiating a noisy spending series tends to amplify measurement error.

The two algorithms implement the same closed-loop pacing idea, illustrated in \Cref{fig:algo_taxonomy}. They differ only in how they measure the error signal and how they translate it into an adjustment of the multiplier.

\begin{figure}[H]
\centering
\begin{tikzpicture}[
  block/.style={rectangle, draw, rounded corners, minimum width=3.0cm,
                minimum height=0.7cm, text centered, font=\small},
  arrow/.style={-{Stealth[length=2mm]}, thick},
  title/.style={font=\small\bfseries, text centered},
]
  \def\colA{0}
  \def\colB{4.5}
  \def\colC{9.0}
  \def\rowA{0}
  \def\rowB{-1.1}
  \def\rowC{-2.2}
  \def\rowD{-3.3}

  \node[title] at (\colA, 0.9) {Q-Learning};
  \node[title] at (\colB, 0.9) {Contextual Bandit};
  \node[title] at (\colC, 0.9) {Pacing};

  \node[block] (q1) at (\colA, \rowA) {Observe state $s_t$};
  \node[block] (q2) at (\colA, \rowB) {$\varepsilon$-greedy pick};
  \node[block] (q3) at (\colA, \rowC) {Submit bid $b_t$};
  \node[block] (q4) at (\colA, \rowD) {Update $Q(s,a)$};

  \node[block] (b1) at (\colB, \rowA) {Observe context $\mathbf{x}_t$};
  \node[block] (b2) at (\colB, \rowB) {UCB or Thompson};
  \node[block] (b3) at (\colB, \rowC) {Submit bid $b_t$};
  \node[block] (b4) at (\colB, \rowD) {Update $\boldsymbol{\theta}_a$};

  \node[block] (p1) at (\colC, \rowA) {Observe $v_t$, budget};
  \node[block] (p2) at (\colC, \rowB) {Shade bid $v_t/\mu_t$};
  \node[block] (p3) at (\colC, \rowC) {Submit bid $b_t$};
  \node[block] (p4) at (\colC, \rowD) {Update dual $\mu_t$};

  \foreach \col in {q, b, p} {
    \draw[arrow] (\col 1) -- (\col 2);
    \draw[arrow] (\col 2) -- (\col 3);
    \draw[arrow] (\col 3) -- (\col 4);
  }

  \draw[arrow, rounded corners=8pt] (q4.west) -- ++(-1.0, 0) |- (q1.west);
  \draw[arrow, rounded corners=8pt] (b4.west) -- ++(-1.0, 0) |- (b1.west);
  \draw[arrow, rounded corners=8pt] (p4.east) -- ++(1.0, 0) |- (p1.east);
\end{tikzpicture}
\caption{The per-round decision loop for each algorithm family. Each column shows the sequence of observation, action selection, bid submission, and parameter update, with a feedback arrow indicating that the cycle repeats.}
\label{fig:algo_taxonomy}
\end{figure}

\section{Equilibrium Benchmarks}
\label{sec:equilibria_main}

Each experiment is paired with a rational-bidder benchmark computed from the same auction environment. These benchmarks serve as yardsticks against which observed algorithmic revenue can be compared. Under constant valuations and risk-neutral bidders, revenue equivalence \citep{Vickrey1961} implies that first-price and second-price auctions generate identical expected revenue. Under affiliated signals, \citet{MilgromWeber1982} show that second-price auctions generate weakly higher expected revenue. When bidders are budget-constrained and pace their bids, the equilibrium concept shifts from Bayes-Nash to pacing equilibrium, and first-price formats acquire uniqueness and efficiency guarantees from \citet{Conitzer2022FPPE}. All derivations and the computed equilibrium bid functions are deferred to Appendix~\ref{ch:equilibria}.

\section{Liquid Welfare and the Price of Anarchy}
\label{sec:welfare}

Experiments~3a and~3b require an efficiency benchmark that accounts for budget constraints. Classical allocative efficiency, defined as the expected value captured by the winning bidders, ignores the fact that a bidder cannot pay more than her budget. Liquid welfare is an alternative measure introduced in the autobidding literature that truncates each bidder's value contribution at her budget. Following \citet{Gaitonde2023} and \citet{Deng2021TowardsEfficient}, the liquid welfare of an allocation $x$ with budgets $B^i$ is
\begin{equation}
W(x) \;=\; \sum_{i = 1}^{n} \min\!\Big\{B^i,\; \sum_{t = 1}^{T} x_{it}\, v_{it}\Big\}.
\label{eq:liquid_welfare}
\end{equation}
The truncation by $B^i$ reflects the fact that value earned beyond the budget is unrealisable, since the agent cannot spend more than its budget allows. The optimal liquid welfare $W^*$ is computed by a linear program that allocates each round to the bidder with the highest valuation subject to per-bidder budget caps. I use the fractional-allocation relaxation of the integer allocation problem, which in the present setting coincides with the integer optimum because the linear program's constraint matrix admits an integer-valued optimal solution whenever one exists.\footnote{The argument follows \citet{Gaitonde2023} and \citet{Deng2021TowardsEfficient}. Both papers use the fractional linear program as the liquid-welfare benchmark and show that it provides a valid upper bound on any achievable allocation, integer or fractional.} The Price of Anarchy is the ratio of optimal to realised liquid welfare,
\begin{equation}
\mathrm{PoA} \;=\; \frac{W^*}{W(x^{\text{obs}})},
\label{eq:poa}
\end{equation}
which is bounded below by~$1$, with smaller values indicating higher efficiency. Theoretical bounds from the autobidding literature suggest $\mathrm{PoA} \leq 2$ for both first-price and second-price formats under cumulative budget constraints \citep{Aggarwal2019, Gaitonde2023}. Experiments~3a and~3b report per-instance Price-of-Anarchy estimates drawn from simulation, not worst-case bounds.

\section{What the Reader Should Take Away}
\label{sec:ch3_summary}

The auction environment is a standard repeated single-item market with a discrete bid grid and a reserve price, and the three algorithm families differ only in how agents turn observations into bids. Q-learning stores a table of state-action values and updates it from rewards. Contextual bandits skip state dynamics and learn a context-to-action map through either optimism (LinUCB) or posterior sampling (Thompson Sampling). Budget-constrained pacing shades a valuation by a multiplier that stands in for the shadow price of the budget, with either multiplicative dual updates or a proportional-integral feedback rule. All three families reduce to simple input-output systems whose outputs can be compared on common metrics in the later chapters. The classical revenue-equivalence and affiliation-based ranking results of \citet{Vickrey1961} and \citet{MilgromWeber1982} provide the rational-bidder benchmark for the unconstrained learning families, and the first-price pacing equilibrium of \citet{Conitzer2022FPPE} together with the Price-of-Anarchy bound of \citet{Aggarwal2019} provide the benchmark for the pacing family.
